\begin{table}[!htbp]
\raggedright
\caption{Item distribution and per-person stationarity.}
\label{tab:sup-stationarity}
\footnotesize
\begingroup
\setlength{\tabcolsep}{4pt}
\renewcommand{\arraystretch}{1.08}
\noindent\textit{Panel A.} Skewness and kurtosis of the core momentary items.\par\vspace{2pt}
\begin{tabular*}{\textwidth}{@{\extracolsep{\fill}}lcc@{}}
\toprule
Measure & Skewness & Kurtosis \\
\midrule
Pain & -0.18 & -0.528 \\
Attention & 1.167 & 0.776 \\
Fear & 0.695 & -0.126 \\
Engagement & -0.165 & -1.055 \\
\bottomrule
\end{tabular*}
\par\vspace{6pt}\noindent\textit{Panel B.} Share of individual series flagged by a linear-trend test and by KPSS.\par\vspace{2pt}
\begin{tabular*}{\textwidth}{@{\extracolsep{\fill}}lcc@{}}
\toprule
Variable & Prop. linear trend & Prop. KPSS non-stationary \\
\midrule
Pain intensity & 0.382 & 0.25 \\
Threat value & 0.412 & 0.279 \\
Attention to pain & 0.338 & 0.235 \\
Goal engagement & 0.221 & 0.147 \\
\bottomrule
\end{tabular*}
\par\vspace{3pt}\noindent\parbox{\textwidth}{\footnotesize\textit{Note.} KPSS = Kwiatkowski-Phillips-Schmidt-Shin test (null is stationarity). A within-person detrended sensitivity model is reported in the robustness battery.}
\endgroup
\end{table}
